\section{Regime-Local Reduced-Order Representation}
\label{sec:regime_local_representation}

RAL-MoE-ROM represents a parametrized flow by a collection of overlapping
regime-local reduced charts.  Each chart contains a local affine POD
representation together with the projected operators required to define the
physics-based part of the reduced dynamics.  This section introduces the
parametrized full-order problem, fixes the pressure gauge used throughout the
paper, constructs the local POD representations, and derives the corresponding
Galerkin and Pressure--Poisson operators.  The learned specialist dynamics are
introduced separately in Section~\ref{sec:regime_local_specialists}.

\subsection{Parametrized full-order model and regime cover}
\label{subsec:parametric_regime_decomposition}

Let $\Omega\subset\mathbb{R}^{2}$ denote the spatial domain and let
$\mu\in\mathcal{M}$ be a physical parameter.  After spatial discretization of
the incompressible Navier--Stokes equations, the velocity field
$\boldsymbol{u}(t;\mu)\in\mathbb{R}^{N_u}$ and pressure field
$p(t;\mu)\in\mathbb{R}^{N_p}$ satisfy the semi-discrete system
\begin{equation}
  \dot{\boldsymbol{u}}
  =
  \boldsymbol{f}_{h}(\mu)
  +\mathcal{L}_{h}(\mu)\boldsymbol{u}
  -\mathcal{N}_{h}
   \bigl(\boldsymbol{u},\boldsymbol{u};\mu\bigr)
  -\nabla_{h}(\mu)p,
  \qquad
  \operatorname{div}_{h}(\mu)\boldsymbol{u}=0,
  \label{eq:semidiscrete_incompressible_system}
\end{equation}
where $\mathcal{L}_{h}$ denotes the discrete linear momentum operator,
$\mathcal{N}_{h}$ the bilinear convective operator, and
$\nabla_h$ and $\operatorname{div}_h$ the discrete gradient and divergence.
The term $\boldsymbol{f}_{h}$ contains forcing and contributions associated
with boundary conditions or lifting functions.  The dependence on $\mu$ is
kept explicit to allow the governing operators and boundary data to vary over
the parameter domain.

For incompressible flow, pressure is determined only up to an additive
constant.  Let $\boldsymbol{w}=(w_1,\ldots,w_{N_p})^{\mathsf T}$ denote the
positive finite-volume quadrature weights.  We use the weighted zero-mean
gauge
\begin{equation}
  \Gamma_p(q)
  =
  q-
  \frac{\boldsymbol{w}^{\mathsf T}q}
       {\boldsymbol{w}^{\mathsf T}\boldsymbol{1}}
  \boldsymbol{1},
  \qquad
  p^{\circ}:=\Gamma_p(p),
  \label{eq:pressure_gauge}
\end{equation}
so that
\begin{equation}
  p^{\circ}\in\mathcal{Q}_{0,h}
  :=
  \left\{
    q\in\mathbb{R}^{N_p}:\boldsymbol{w}^{\mathsf T}q=0
  \right\}.
  \label{eq:discrete_zero_mean_pressure_space}
\end{equation}
Because the momentum equation depends only on the pressure gradient,
$p$ in~\eqref{eq:semidiscrete_incompressible_system} can henceforth be
replaced by its gauge-fixed representative $p^{\circ}$.  The full-order state
is therefore written as
\begin{equation}
  \boldsymbol{x}(t;\mu)
  =
  \bigl(\boldsymbol{u}(t;\mu),p^{\circ}(t;\mu)\bigr).
  \label{eq:full_state_definition}
\end{equation}

The parameter domain is represented by an overlapping cover
\begin{equation}
  \mathcal{M}
  \subseteq
  \bigcup_{r\in\mathcal{R}}\mathcal{M}_r,
  \qquad
  \mathcal{R}=\{S,H,P\},
  \label{eq:regime_cover}
\end{equation}
where $S$, $H$, and $P$ denote the steady, Hopf-onset, and periodic regimes,
respectively.  The adjacent overlap regions are
\begin{equation}
  \mathcal{O}_{SH}=\mathcal{M}_{S}\cap\mathcal{M}_{H},
  \qquad
  \mathcal{O}_{HP}=\mathcal{M}_{H}\cap\mathcal{M}_{P}.
  \label{eq:adjacent_overlaps}
\end{equation}
The cover is not required to be disjoint: near a regime transition, two
neighboring reduced descriptions may both provide useful local
approximations.  The corresponding selection and coupling of local models is
deferred to Section~\ref{sec:hierarchical_ral_moe}.

\subsection{Regime-local weighted POD representation}
\label{subsec:local_pod_and_gauge}

The finite-volume quadrature induces the discrete inner products
\begin{equation}
  \langle \boldsymbol{v},\boldsymbol{z}\rangle_{M_u}
  =\boldsymbol{v}^{\mathsf T}M_u\boldsymbol{z},
  \qquad
  \langle q,s\rangle_{M_p}
  =q^{\mathsf T}M_p s,
  \label{eq:weighted_inner_products}
\end{equation}
with associated norms $\|\cdot\|_{M_u}$ and $\|\cdot\|_{M_p}$.  For the
collocated two-dimensional finite-volume discretization considered here,
\begin{equation}
  M_p=\operatorname{diag}(w_1,\ldots,w_{N_p}),
  \qquad
  M_u=\operatorname{diag}(M_p,M_p).
  \label{eq:velocity_pressure_mass_matrices}
\end{equation}

For each regime $r\in\mathcal{R}$, let
\begin{equation}
  \mathcal{S}_r
  =
  \left\{
    \bigl(\boldsymbol{u}^{(r)}_s,p^{(r),\circ}_s\bigr)
  \right\}_{s=1}^{N_r},
  \qquad
  p^{(r),\circ}_s=\Gamma_p\bigl(p^{(r)}_s\bigr),
  \label{eq:regime_snapshot_ensemble}
\end{equation}
be the snapshot ensemble used to construct the local representation.  The precise trajectory partition employed in the numerical experiments is
stated with the common evaluation protocol; it is not part of the definition
of a local chart.

The regime-local means are
\begin{equation}
  \overline{\boldsymbol{u}}^{(r)}
  =\frac{1}{N_r}\sum_{s=1}^{N_r}\boldsymbol{u}^{(r)}_s,
  \qquad
  \overline{p}^{(r),\circ}
  =\frac{1}{N_r}\sum_{s=1}^{N_r}p^{(r),\circ}_s,
  \label{eq:local_training_means}
\end{equation}
and the centered snapshot matrices are
\begin{align}
  X_u^{(r)}
  &=
  \left[
    \boldsymbol{u}^{(r)}_1-\overline{\boldsymbol{u}}^{(r)}
    \ \cdots\
    \boldsymbol{u}^{(r)}_{N_r}-\overline{\boldsymbol{u}}^{(r)}
  \right],
  \label{eq:centered_velocity_snapshots}\\
  X_p^{(r)}
  &=
  \left[
    p^{(r),\circ}_1-\overline{p}^{(r),\circ}
    \ \cdots\
    p^{(r),\circ}_{N_r}-\overline{p}^{(r),\circ}
  \right].
  \label{eq:centered_pressure_snapshots}
\end{align}

The local velocity and pressure POD bases are the weighted orthonormal
minimizers
\begin{align}
  \Phi_u^{(r)}
  &\in
  \underset{\Psi^{\mathsf T}M_u\Psi=I_{d_u}}
  {\operatorname{arg\,min}}
  \left\|
    X_u^{(r)}-\Psi\Psi^{\mathsf T}M_uX_u^{(r)}
  \right\|_{M_u,F}^{2},
  \label{eq:velocity_pod_minimization}\\
  \Phi_p^{(r)}
  &\in
  \underset{\Psi^{\mathsf T}M_p\Psi=I_{d_p}}
  {\operatorname{arg\,min}}
  \left\|
    X_p^{(r)}-\Psi\Psi^{\mathsf T}M_pX_p^{(r)}
  \right\|_{M_p,F}^{2},
  \label{eq:pressure_pod_minimization}
\end{align}
where $\|Y\|_{M,F}^{2}=\operatorname{tr}(Y^{\mathsf T}MY)$.  Equivalently,
if
\begin{align}
  M_u^{1/2}X_u^{(r)}
  &=U_u^{(r)}\Sigma_u^{(r)}{V_u^{(r)}}^{\mathsf T},\\
  M_p^{1/2}X_p^{(r)}
  &=U_p^{(r)}\Sigma_p^{(r)}{V_p^{(r)}}^{\mathsf T},
\end{align}
then
\begin{equation}
  \Phi_u^{(r)}
  =M_u^{-1/2}U_u^{(r)}(:,1\!:\!d_u),
  \qquad
  \Phi_p^{(r)}
  =M_p^{-1/2}U_p^{(r)}(:,1\!:\!d_p),
  \label{eq:weighted_pod_bases}
\end{equation}
with
\begin{equation}
  {\Phi_u^{(r)}}^{\mathsf T}M_u\Phi_u^{(r)}=I_{d_u},
  \qquad
  {\Phi_p^{(r)}}^{\mathsf T}M_p\Phi_p^{(r)}=I_{d_p}.
  \label{eq:weighted_pod_orthogonality}
\end{equation}
Since the pressure snapshots and their mean belong to
$\mathcal{Q}_{0,h}$, the pressure modes satisfy
$\boldsymbol{w}^{\mathsf T}\Phi_p^{(r)}=\boldsymbol{0}^{\mathsf T}$.

The local reduced coordinates are obtained by weighted orthogonal projection,
\begin{align}
  \boldsymbol{a}^{(r)}
  &=
  {\Phi_u^{(r)}}^{\mathsf T}M_u
  \bigl(\boldsymbol{u}-\overline{\boldsymbol{u}}^{(r)}\bigr),
  \label{eq:velocity_encode}\\
  \boldsymbol{b}^{(r)}
  &=
  {\Phi_p^{(r)}}^{\mathsf T}M_p
  \bigl(\Gamma_p(p)-\overline{p}^{(r),\circ}\bigr),
  \label{eq:pressure_encode}
\end{align}
and the corresponding affine reconstructions are
\begin{align}
  \mathcal{D}^{u}_{r}\bigl(\boldsymbol{a}^{(r)}\bigr)
  &=
  \overline{\boldsymbol{u}}^{(r)}
  +\Phi_u^{(r)}\boldsymbol{a}^{(r)},
  \label{eq:velocity_decode}\\
  \mathcal{D}^{p}_{r}\bigl(\boldsymbol{b}^{(r)}\bigr)
  &=
  \overline{p}^{(r),\circ}
  +\Phi_p^{(r)}\boldsymbol{b}^{(r)}.
  \label{eq:pressure_decode}
\end{align}
For the learned components introduced later, the modal coordinates are
standardized locally,
\begin{equation}
  \widetilde{\boldsymbol{a}}^{(r)}
  =
  \bigl(\boldsymbol{a}^{(r)}-\boldsymbol{\mu}_{a}^{(r)}\bigr)
  \oslash\boldsymbol{\sigma}_{a}^{(r)},
  \qquad
  \widetilde{\boldsymbol{b}}^{(r)}
  =
  \bigl(\boldsymbol{b}^{(r)}-\boldsymbol{\mu}_{b}^{(r)}\bigr)
  \oslash\boldsymbol{\sigma}_{b}^{(r)},
  \label{eq:local_scalers}
\end{equation}
where $\oslash$ denotes componentwise division.  These scalings are local to
each chart and are not used in the projected physical operators.

\subsection{Chart-local POD--Galerkin and Pressure--Poisson operators}
\label{subsec:local_galerkin_pressure_poisson}

Substituting the affine POD reconstructions into
\eqref{eq:semidiscrete_incompressible_system} and projecting the momentum
equation onto the local velocity space gives
\begin{equation}
  \dot{\boldsymbol{a}}^{(r)}
  =
  \boldsymbol{c}_r(\mu)
  +A_r(\mu)\boldsymbol{a}^{(r)}
  +H_r(\mu):
   \bigl(\boldsymbol{a}^{(r)}\otimes\boldsymbol{a}^{(r)}\bigr)
  +C_r^p(\mu)\boldsymbol{b}^{(r)}.
  \label{eq:projected_local_momentum_general_mass}
\end{equation}
The reduced operators are defined by
\begin{align}
  [\boldsymbol{c}_r(\mu)]_i
  &=
  \left\langle
    \boldsymbol{\phi}^{u,(r)}_i,
    \boldsymbol{f}_h
    +\mathcal{L}_h\overline{\boldsymbol{u}}^{(r)}
    -\mathcal{N}_h
      (\overline{\boldsymbol{u}}^{(r)},\overline{\boldsymbol{u}}^{(r)};\mu)
    -\nabla_h\overline{p}^{(r),\circ}
  \right\rangle_{M_u},
  \label{eq:galerkin_constant_definition}\\
  [A_r(\mu)]_{ij}
  &=
  \left\langle
    \boldsymbol{\phi}^{u,(r)}_i,
    \mathcal{L}_h\boldsymbol{\phi}^{u,(r)}_j
    -\mathcal{N}_h
      (\overline{\boldsymbol{u}}^{(r)},\boldsymbol{\phi}^{u,(r)}_j;\mu)
    -\mathcal{N}_h
      (\boldsymbol{\phi}^{u,(r)}_j,\overline{\boldsymbol{u}}^{(r)};\mu)
  \right\rangle_{M_u},
  \label{eq:galerkin_linear_definition}\\
  [H_r(\mu)]_{ijk}
  &=
  -\left\langle
    \boldsymbol{\phi}^{u,(r)}_i,
    \mathcal{N}_h
      (\boldsymbol{\phi}^{u,(r)}_j,\boldsymbol{\phi}^{u,(r)}_k;\mu)
  \right\rangle_{M_u},
  \label{eq:galerkin_quadratic_definition}\\
  [C_r^p(\mu)]_{i\ell}
  &=
  -\left\langle
    \boldsymbol{\phi}^{u,(r)}_i,
    \nabla_h(\mu)\phi^{p,(r)}_\ell
  \right\rangle_{M_u}.
  \label{eq:galerkin_pressure_coupling_definition}
\end{align}
Parameter dependence has been suppressed inside some operator arguments for
readability.  The resulting local Galerkin vector field is
\begin{equation}
  \mathscr{F}^{\mathrm{G}}_r
  (\boldsymbol{a},\boldsymbol{b};\mu)
  :=
  \boldsymbol{c}_r(\mu)
  +A_r(\mu)\boldsymbol{a}
  +H_r(\mu):(\boldsymbol{a}\otimes\boldsymbol{a})
  +C_r^p(\mu)\boldsymbol{b}.
  \label{eq:chart_local_galerkin_rhs_definition}
\end{equation}

Pressure is reconstructed algebraically from the incompressibility constraint
and momentum balance.  At the continuous level, the gauge-fixed pressure
satisfies a Pressure--Poisson problem of the form
\begin{equation}
  -\Delta p^{\circ}
  =
  \nabla\!\cdot
  \left[(\boldsymbol{u}\!\cdot\!\nabla)\boldsymbol{u}
        -\boldsymbol{f}\right],
  \qquad
  \int_{\Omega}p^{\circ}\,\mathrm{d}\Omega=0,
  \label{eq:pressure_poisson_strong_form}
\end{equation}
with boundary conditions inherited from the momentum equation.  Projecting
the corresponding discrete Pressure--Poisson equation onto the local pressure
space gives
\begin{equation}
  K_r^p(\mu)\boldsymbol{b}^{PP,(r)}
  =
  \boldsymbol{d}_r(\mu)
  +B_r(\mu)\boldsymbol{a}^{(r)}
  +Q_r(\mu):
   \bigl(\boldsymbol{a}^{(r)}\otimes\boldsymbol{a}^{(r)}\bigr).
  \label{eq:reduced_pressure_poisson_system}
\end{equation}
Here $K_r^p$, $\boldsymbol{d}_r$, $B_r$, and $Q_r$ are obtained by testing
the discrete Pressure--Poisson operator and its constant, linear, and
quadratic velocity contributions against the local pressure modes.  The
pressure gauge removes the constant nullspace, so that the reduced algebraic
map is well defined under the standard coercivity assumptions:
\begin{equation}
  \mathscr{Q}^{PP}_r(\boldsymbol{a};\mu)
  =
  [K_r^p(\mu)]^{-1}
  \left[
    \boldsymbol{d}_r(\mu)
    +B_r(\mu)\boldsymbol{a}
    +Q_r(\mu):(\boldsymbol{a}\otimes\boldsymbol{a})
  \right].
  \label{eq:chart_local_pressure_poisson_map}
\end{equation}
The map $\mathscr{Q}^{PP}_r$ supplies the physics-based pressure relation used
by the regime-local specialist in Section~\ref{sec:regime_local_specialists}.

\subsection{Complete regime-local chart and coordinate independence}
\label{subsec:complete_regime_local_rom_chart}

The regime-local reduced representation is collected as
\begin{equation}
  \mathcal{C}^{\mathrm{ROM}}_r
  =
  \left(
    \mathcal{M}_r,
    \Gamma_p,
    \overline{\boldsymbol{u}}^{(r)},
    \Phi_u^{(r)},
    \overline{p}^{(r),\circ},
    \Phi_p^{(r)},
    \boldsymbol{\mu}_a^{(r)},
    \boldsymbol{\sigma}_a^{(r)},
    \boldsymbol{\mu}_b^{(r)},
    \boldsymbol{\sigma}_b^{(r)},
    \mathscr{F}^{\mathrm{G}}_r,
    \mathscr{Q}^{PP}_r
  \right).
  \label{eq:local_chart_contract}
\end{equation}
Thus, a chart consists of a parameter region, a gauge-consistent affine POD
representation, local coefficient scalings, and the projected operators used
to define its physics-based reduced dynamics.

Crucially, different charts are independent coordinate systems.  Even when
two charts use the same dimensions $d_u$ and $d_p$, their affine means and POD
bases are obtained from different portions of the solution manifold.  Hence
there is, in general, no canonical componentwise identification between
$\boldsymbol{a}^{(r)}$ and $\boldsymbol{a}^{(k)}$, or between
$\boldsymbol{b}^{(r)}$ and $\boldsymbol{b}^{(k)}$.  A physical state in an
overlap can instead be encoded independently by each local projection
\eqref{eq:velocity_encode}--\eqref{eq:pressure_encode}.  This observation is
central to the hierarchical construction developed later: neighboring
specialists evolve in their own reduced coordinates, while any soft coupling
is performed only after reconstruction in the common physical space.
